\documentclass[a4paper,10pt]{ltjsarticle}
\usepackage[margin=19mm,top=17mm,bottom=18mm]{geometry}
\usepackage{amsmath,amssymb,booktabs}
\usepackage{luatexja-fontspec}
%\setmainjfont{HaranoAjiMincho}
%\setsansjfont{HaranoAjiGothic}
\usepackage[colorlinks=true,linkcolor=blue]{hyperref}
\newcommand{\PP}{\mathbb P}
\newcommand{\ZZ}{\mathbb Z}
\newcommand{\CC}{\mathbb C}
\newcommand{\Pic}{\operatorname{Pic}}
\newcommand{\Aut}{\operatorname{Aut}}
\newcommand{\im}{\operatorname{Im}}
\newcommand{\LL}{\mathcal L}
\newcommand{\Pn}{\mathcal P}
\newcommand{\I}{\mathcal I}
\newcommand{\head}[2]{\par\medskip\noindent{\large\bfseries #1}\hfill{\small #2}\par\medskip}
\newcommand{\cue}[1]{\par\smallskip\noindent{\small\sffamily［板書・進行］#1}\par\smallskip}
\setlength{\parindent}{1em}
\setlength{\parskip}{3pt}
\linespread{1.08}
\pagestyle{plain}
\begin{document}
\begin{center}
{\Large\bfseries Del Pezzo 曲面の族に付随するモノドロミー表現}\par\smallskip
{\large Del Pezzo 曲面の族に付随するモノドロミー表現}\par\smallskip
島田伊知朗\quad 代数セミナー講演原稿案 第2版\quad 2026年10月
\end{center}
%\noindent{\small 実質80分，質疑等10分．導入30分，Corollary 3.2 と全射性40分，Proposition 4.4 への応用10分を目安とする．板書を挟んで使う講演原稿である．以下，有限体上の計算例を除き，基礎体は $\CC$ とする．}
\head{1\quad del Pezzo 曲面と一般の位置}{0〜15分}
滑らかな3次曲面には27本の直線がある．曲面を動かして元へ戻ると，これらの直線には置換が生じる．交点配置から許される Weyl 群の対称性は，実はすべてこのような変形で実現する．本講演では次数1の del Pezzo 曲面を主な対象とし，240本の lines のモノドロミーが $W(E_8)$ に全射となることを証明する．

\noindent\textbf{定義．}\quad del Pezzo 曲面とは，反標準因子 $-K_X$ が ample である滑らかな射影曲面 $X$ である．$\alpha=[-K_X]$ と書き，$d=\alpha^2$ を次数と呼ぶ．ample とは，ある正の整数倍が射影空間への埋め込みを与えるという性質である．$\PP^2$ は次数9，$\PP^1\times\PP^1$ は次数8の例である．

blow up は平面の一点をその点での方向を表す $\PP^1$ に置き換える操作であり，この例外曲線の自己交点は $-1$ である．相異なる順序付き $n$ 点 $p=(p_1,\ldots,p_n)$ における blow up を $\beta_p:Y_p\to\PP^2$ と書く．

\noindent\textbf{定義 3.4（一般の位置）．}\quad $1\leq n\leq8$ とする．次の条件をすべて満たす相異なる順序付き $n$ 点の集合を $\Pn_n\subset(\PP^2)^n$ とする．
\begin{enumerate}
\setlength{\itemsep}{0pt}\setlength{\parskip}{0pt}
\item どの3点も一直線上にない．
\item どの6点も一つの2次曲線上にない．
\item どの8点についても，そのうちの7点を通り，残りの1点に二重点を持つ3次曲線は存在しない．
\end{enumerate}
\noindent ここで2次曲線と3次曲線は可約なものも含む．(3) はどの点を二重点に選んでも成立する条件であり，二重点には重複度が2以上の場合を含める．点数が不足する条件は空条件である．$\Pn_n$ は空でない Zariski 開集合であり，特に既約である．

\noindent\textbf{定理 3.5．}\quad $Y_p$ が次数 $d=9-n$ の del Pezzo 曲面であることと，$p\in\Pn_n$ であることは同値である．

基本的な分類によれば，del Pezzo 曲面は $\PP^1\times\PP^1$ または平面の一般の位置の $0\leq n\leq8$ 点の blow up である．以下では主に $1\leq d\leq6$ を扱う．一般の位置が必要な理由は，例えば3点を通る直線の狭義変換 $C$ では $\alpha\cdot C=3-3=0$ となり，$\alpha$ の ample 性に反することから見える．

\noindent\textbf{定義．}\quad 滑らかな有理曲線 $l\subset X$ で $\alpha\cdot l=1$ となるものを line と呼ぶ．随伴公式より $-2=(K_X+l)\cdot l=-1+l^2$ なので，$l^2=-1$ である．
\[
\begin{array}{c|rrrrrr}
d&6&5&4&3&2&1\\\hline
\text{lines の本数}&6&10&16&27&56&240
\end{array}
\]
次数3では反標準埋め込みで実際に直線となる．平面の6点の blow up で見ると，例外曲線6本，2点を結ぶ直線の狭義変換15本，5点を通る2次曲線の狭義変換6本で，計27本である．
\cue{blow up と27本の三種類を図示する．定義3.4の条件は，後半の有限体上の例でも確認すべき条件となる．}
\newpage

\head{2\quad Picard 格子からルート格子へ}{15〜30分}
曲線を交点数で調べるために，Picard 群を使う．この曲面では $\Pic(X)\cong H^2(X,\ZZ)$ であり，交点形式を入れたものを Picard 格子と呼ぶ．平面の直線の引き戻しを $h$，例外曲線の類を $e_1,\ldots,e_n$ とすると，
\[
\Pic(X)=\ZZ h\oplus\bigoplus_{i=1}^{n}\ZZ e_i,\qquad
h^2=1,\quad e_i^2=-1,\quad h\cdot e_i=0,\quad e_i\cdot e_j=0\ (i\ne j),
\]
\[
\alpha=3h-e_1-\cdots-e_n,\qquad \alpha^2=9-n=d.
\]
ここで $\alpha$ に直交する部分を取り出す．
\[
R(X):=\alpha^\perp\subset\Pic(X).
\]
Hodge の指数定理からこれは負定値である．具体的に，
\[
r_i=e_i-e_{i+1}\ (1\leq i\leq n-1),\qquad
r_n=h-e_1-e_2-e_3
\]
をとると，どれも $\alpha$ と直交し，自己交点は $-2$ である．これらは $R(X)$ の基底になる．異なる二つの交点数が1のとき辺で結ぶと，Dynkin 図形が現れる．
\[
\begin{array}{c|rrrrrr}
d&6&5&4&3&2&1\\\hline
R(X)\text{ の型}&A_1+A_2&A_4&D_5&E_6&E_7&E_8
\end{array}
\]
今日は負定値の符号で統一する．ルートとは $r^2=-2$ のベクトルである．ルート $r$ に関する反射は
\[
s_r(x)=x+(x\cdot r)r
\]
であり，$r$ を $-r$ にし，$r^\perp$ を固定する．これらの反射で生成される群が Weyl 群 $W(R(X))$ である．例えば $r=e_1-e_2$ の反射は $e_1,e_2$ を交換する．$h-e_1-e_2-e_3$ の反射は平面の標準的な2次 Cremona 変換に対応する基底の変換である．

ここで使う重要な事実が論文の Proposition 3.1 である．
\[
\boxed{\ O(\Pic(X),\alpha)\ \cong\ W(R(X)).\ }
\]
左辺は交点形式と反標準類を保つ格子の対称性である．ルート格子の等長変換には Dynkin 図形の対称性もあり得るが，$\alpha$ を固定して Picard 格子全体に延長できるものはちょうど Weyl 群になる．一般の次数の証明には判別群を用いる．次数1では $\alpha^2=1$ なので
\[
\Pic(X)=\ZZ\alpha\ \perp\ R(X),\qquad R(X)\cong E_8,
\]
と直交分解し，$E_8$ の Dynkin 図形に非自明な対称性がないことから直ちに分かる．

さらに次数1では，lines とルートが直接対応する．
\[
L(X)\longrightarrow\Delta(E_8),\qquad l\longmapsto[l]-\alpha.
\]
実際 $([l]-\alpha)\cdot\alpha=0$，$([l]-\alpha)^2=-2$ である．逆にルート $r$ から $\alpha+r$ が line の類になる．ここでは「$\lambda^2=-1$, $\alpha\cdot\lambda=1$ を満たす類は一意な line で表される」という del Pezzo 曲面の基本事実を用いる．したがって240本という数は $E_8$ の240個のルートの数である．
\cue{$E_8$ の図形は $r_1,\ldots,r_7$ の鎖の $r_3$ に $r_8$ を付ける．ここまでで30分．Weyl 群の元がすべて固定した曲面の自己同型になる，という主張ではないことを一言添える．}
\newpage
\head{3\quad Corollary 3.2 と全射性の判定法}{30〜45分}
ここからが今日の中心である．一本の line の代わりに，互いに交わらない $n=9-d$ 本の lines を，順序も付けて選ぶ．その集合を
\[
L^{[n]}(X)=\{(l_1,\ldots,l_n)\mid l_i\text{ は line},\ l_i\cdot l_j=0\ (i\ne j)\}
\]
と書く．この $n$ 本を同時に縮約すると $\PP^2$ になり，$X$ を平面のどの点の blow up と見るか，という表示が得られる．縮約後の平面の座標までは選んでいない．

\noindent\textbf{Corollary 3.2.}\quad
$W(R(X))\cong O(\Pic(X),\alpha)$ の $L^{[n]}(X)$ への作用は自由かつ推移的である．

つまり，二つの順序付き配置を指定すると，一方を他方に移す Weyl 群の元が，ただ一つ存在する．一本の line 上の作用が推移的，というだけよりもずっと強い主張である．

この主張を基底の言葉で説明する．$\lambda=(l_1,\ldots,l_n)$ を縮約する写像を $\beta_\lambda$ とし，平面の直線の引き戻しを $h_\lambda$ とすると，
\[
3h_\lambda=\alpha+l_1+\cdots+l_n.
\]
したがって配置から $h_\lambda$ も決まり，$h_\lambda,l_1,\ldots,l_n$ はいつもの交点行列 $\operatorname{diag}(1,-1,\ldots,-1)$ を持つ基底になる．別の配置 $\mu$ についても同じである．二つの基底を対応させれば，$\alpha$ を保つ格子の等長変換が得られる．Proposition 3.1 により，それは Weyl 群の元である．これが推移性である．また各 $l_i$ と $\alpha$ を固定すれば $h_\lambda$ も固定するので，基底全体を固定する．これが自由性である．

次に，del Pezzo 曲面の滑らかな族 $f:\mathcal X\to U$ を考える．$U$ は滑らかで既約とし，基点を $b$ とする．$U$ のループに沿って曲面のコホモロジーを運ぶと，交点形式と反標準類を保つ変換が得られる．これがモノドロミー
\[
\Phi:\pi_1(U,b)\longrightarrow O(\Pic(X_b),\alpha_b)=W(R(X_b))
\]
である．問題は $\Gamma:=\im\Phi$ が $W:=W(R(X_b))$ 全体になるかどうかである．

各曲面で順序付きの配置を選ぶ空間を作り，有限エタール被覆
\[
\LL^{[n]}\longrightarrow U,\qquad (\LL^{[n]})_u=L^{[n]}(X_u)
\]
を考える．被覆空間の基本的な事実として，全空間の連結成分は，基点上のファイバーにおけるモノドロミーの軌道と一対一に対応する．Corollary 3.2 によってファイバーは $W$ の自由推移的な集合であるから，$\Gamma$ の軌道の個数は指数 $[W:\Gamma]$ である．したがって
\[
\boxed{\ [W:\im\Phi]=\#\pi_0(\LL^{[n]}).\ }\qquad\text{（Corollary 3.3）}
\]
特に，この被覆の全空間が連結であることと，モノドロミーが全射であることは同値である．群の元を一つずつループで実現する代わりに，配置を付けた曲面の空間の連結性を調べればよいのである．
\cue{被覆の二つの点が道で結べることと，下のループによって互いに移ることを図示する．「順序付き」が不可欠で，順序を忘れると安定化群が残ることを強調する．}
\newpage
\head{4\quad 平面の点配置を使って連結成分を数える}{45〜60分}
この判定法を使うために，曲面の族と平面の点配置を結ぶ．$\Pn_n$ を，先ほどの一般の位置の条件を満たす順序付き $n$ 点 $p=(p_1,\ldots,p_n)$ の空間とする．これは $(\PP^2)^n$ の空でない Zariski 開集合なので既約であり，特に連結である．$Y_p$ をその点配置の blow up と書く．

両方の表示を一緒に記録する空間を作る．
\[
\I=\{(u,p,g)\mid u\in U,\ p\in\Pn_n,\ g:X_u\xrightarrow{\sim}Y_p\}.
\]
$Y_p$ の順序付き例外曲線を $g$ で引き戻すと，$X_u$ の順序付き配置が得られる．したがって二つの写像
\[
\LL^{[n]}\ \longleftarrow\ \I\ \longrightarrow\ \Pn_n
\]
がある．左の写像のファイバーは，$n$ 本を縮約した平面と，あらかじめ固定した $\PP^2$ との同型の空間である．これは $\operatorname{PGL}(3,\CC)$ と同型で連結である．この写像は平面の座標を選ぶ束なので，$\I$ と $\LL^{[n]}$ の連結成分数は一致する．残る仕事は右の写像のファイバーを調べることである．

まず次数3で適用する．$U\subset|\mathcal O_{\PP^3}(3)|$ を滑らかな3次曲面の空間とし，$n=6$ とする．$p$ を固定すると，$Y_p$ の反標準線形系は3次曲面への埋め込みを与える．それを固定した $\PP^3$ に置く自由度は
\[
\operatorname{Isom}(|-K_{Y_p}|^\vee,\PP^3)\cong\operatorname{PGL}(4,\CC)
\]
である．基底 $\Pn_6$ もこのファイバーも連結なので $\I$ は連結である．よって $\LL^{[6]}$ は連結，従って $\Phi$ は $W(E_6)$ に全射である．これが27本の直線のモノドロミーの証明である（Proposition 3.6）．

いよいよ次数1である．$|-K_X|$ 自体は埋め込みを与えないが，$|-2K_X|$ は2次錐への2重被覆
\[
X\longrightarrow Q\subset\PP^3
\]
を与える．$Q$ の頂点を $V$ とすると，分岐は $B\cup\{V\}$ で，$B$ は $|\mathcal O_Q(3)|$ の滑らかな曲線である．固定した $Q$ 上で，このような $B$ を動かすパラメータ空間 $U\subset|\mathcal O_Q(3)|\cong\PP^{15}$ を考える（論文 \S3.5 の族）．

$p\in\Pn_8$ を固定し，$Y_p\to Q_p$ をその双反標準モデルとする．同型 $g:X_u\to Y_p$ は同型 $\bar g:Q\to Q_p$ を誘導する．逆に $\gamma:Q\to Q_p$ を与えると，分岐曲線の逆像によって $u$ が決まり，$\gamma$ の持ち上げはちょうど二つある．二つは2重被覆の被覆変換だけ異なる．

$\operatorname{Isom}(Q,Q_p)$ は連結な滑らかな代数多様体である．実際，2次錐の自己同型群は連結であり，これは底の $\PP^1$ の座標変換，錐の方向の非零スカラー倍，および2次形式を加える変換から分かる．その上の二つの持ち上げを選ぶ空間は高々二つの連結成分を持つ．この構成を $\Pn_8$ 上で局所的に行うと，$\I$ も高々二成分である．従って
\[
\boxed{\ [W(E_8):\Gamma]\leq2.\ }
\]
ここでは，持ち上げが二つあることから「連結」と断定はできない．最後に残った指数2の可能性を，一つの局所モノドロミーで排除する．
\newpage
\head{5\quad 局所モノドロミーと有限単純群}{60〜70分}
分岐曲線を，頂点 $V$ を避け，一つだけ節点を持つ曲線 $B_q$ に退化させる．対応する2重被覆 $X_q$ は普通の二重点を一つ持ち，局所的な平滑化は
\[
z_1^2+z_2^2+z_3^2=t
\]
である．$t\to0$ でつぶれる実2次元の球面の類を消滅サイクル $v$ と呼ぶ．Poincar\'e 双対で Picard 格子の元と見れば，$v^2=-2$，$\alpha_b\cdot v=0$ であり，$v$ はルートである．特異値を一周するループのモノドロミーは Picard--Lefschetz の公式
\[
x\longmapsto x+(x\cdot v)v=s_v(x)
\]
で与えられる．従って $\Gamma=\im\Phi$ は反射 $s_v$ を含み，$\det(s_v)=-1$ である．

\noindent\textbf{$W(E_8)$ と有限単純群．}\quad $W=W(E_8)$，$W^+=\ker(\det:W\to\{\pm1\})$ と置く．中心には $Z=\{\pm\mathrm{id}\}$ があり，階数が8なので $-\mathrm{id}\in W^+$ である．論文で使う群論的事実は
\[
G:=W^+/Z\quad\text{が非可換有限単純群である}
\]
ことである．その関係は二つの完全列で表される．
\[
1\longrightarrow Z\longrightarrow W^+\longrightarrow G\longrightarrow1,
\qquad
1\longrightarrow G\longrightarrow W/Z\xrightarrow{\det}\{\pm1\}\longrightarrow1.
\]
論文の記法では $W^+=2.G$，$W/Z=G.2$，$W=2.G.2$ である．ここで点は群の拡大を表し，直積を意味しない．位数は
\[
|W|=696729600,\quad |W^+|=|W/Z|=348364800,\quad |G|=174182400
\]
となる．240個のルートに作用するのが $W$，符号を忘れた120組のルート対に作用するのが $W/Z$ である．この中心の符号は，幾何的には Bertini involution に対応する．

\noindent\textbf{Proposition 3.8 の証明の結び．}\quad 既に $[W:\Gamma]\leq2$ が分かっている．指数が2と仮定し，$H=\Gamma\cap W^+$ と置く．$s_v\in\Gamma$ の行列式が $-1$ なので $[\Gamma:H]=2$，従って $[W^+:H]=2$ である．もし $-\mathrm{id}\in\Gamma$ なら $Z\subset H$ となり，$H/Z$ は単純群 $G$ の指数2の部分群になる．指数2の部分群は正規なので，これは単純性に反する．

従って $\Gamma\cap Z=\{1\}$ である．また $|\Gamma|=|W/Z|$ なので，射影は同型 $\Gamma\xrightarrow{\sim}W/Z$ を与える．$W/Z$ はルート対 $\{\pm r\}$ 全体に推移的に作用する．ゆえに任意のルート $r$ に対して，ある $g\in\Gamma$ が $v$ を $\pm r$ に移す．反射は符号に依存しないため
\[
gs_vg^{-1}=s_{g(v)}=s_r\in\Gamma.
\]
すべてのルート反射が $\Gamma$ に属し，$W$ はそれらで生成されるので $\Gamma=W$ である．これは指数2に矛盾する．以上より
\[
\boxed{\ \Phi:\pi_1(U,b)\twoheadrightarrow W(E_8).\ }
\]
\cue{ここでは論文の有限単純群を使う証明を説明した．短縮する場合には「指数2なら $\Gamma\triangleleft W$，すべてのルート反射は共役，従って一つ含めばすべて含む」としてもよい．}
\newpage
\head{6\quad 応用としての240本の lines の交わり}{70〜75分}
モノドロミーの全射性は，幾何的な性質を確かめるための計算量を減らす道具になる．その例が次の命題である．

\noindent\textbf{命題 4.4．}\quad 一般の次数1の del Pezzo 曲面 $X_b$ において，240本の lines の合併は，特異点として通常二重点のみを持つ．

ここで「一般の」はパラメータ空間のある空でない Zariski 開集合上という意味である．すべての次数1の曲面についての主張ではない．各 line は滑らかなので，示すべきことは，異なる2本が接しないこと，および異なる3本が一点を共有しないことである．前者は2本の横断性，後者は三重点の排除である．

\noindent\textbf{二つ組の軌道．}\quad $r_i=[l_i]-\alpha$ とすると，
\[
l_i\cdot l_j=1+r_i\cdot r_j.
\]
$E_8$ のルートの組合せを調べると，異なる2本の交点数は $m=0,1,2,3$ のいずれかである．$W(E_8)$ の二つ組への作用にはちょうど4軌道があり，交点数 $m$ が軌道を区別する（\S4.2）．全射性により，交点数 $m$ の二つ組を付けた族 $\LL^{\{2\}}_m$ は既約である（Corollary 4.3）．

$m=0$ なら交わらず，$m=1$ なら局所交点数の正値性により一点で横断的に交わる．従って $m=2,3$ のそれぞれで，異なる $m$ 点で交わる一例を見つければよい．横断性は開条件なので，その軌道の一般の配置でも同じ性質が成り立つ．

\noindent\textbf{三つ組の軌道．}\quad 三つの交点数を小さい順に並べたものを $t$ と書く．三つ組には12軌道があり，$t$ によって区別される．その一覧は次である．
\[
\begin{gathered}
[0,0,0],\ [0,0,1],\ [0,0,2],\ [0,1,1],\ [0,1,2],\ [0,2,2],\ [0,2,3],\\
[1,1,1],\ [1,1,2],\ [1,1,3],\ [1,2,2],\ [2,2,2].
\end{gathered}
\]
交点数に0を含む場合には，三つの曲線が共通点を持つことはない．従って調べるべきものは下段の5軌道だけである．各軌道で $l_1\cap l_2\cap l_3=\varnothing$ となる一例を与えればよい．共通点を持たないことも開条件である．

このように命題4.4は，二つ組について2種類，三つ組について5種類の存在確認に帰着する．被覆は有限なので，各既約成分の悪い閉集合の像も底の真の閉集合である．それら有限個の像を除けば，その曲面上のすべての組が同時に良い性質を持つ．これで「一般の配置」から「一般の曲面の240本全部」への移行もできる．
\cue{軌道数4と12はルート配置の有限計算による結果として述べる．全射性がなければ，一つの Weyl 群軌道が複数のモノドロミー軌道に分かれる可能性があり，一例の確認だけでは足りない．}
\newpage
\head{7\quad Proposition 4.4 の計算例と証明の概要}{75〜80分}
残るのは，前頁の7種類の条件を満たす例を見つけることである．論文では $\mathbb F_{19}$ 上の平面の次の8点を用いる．
\[
\begin{array}{llll}
p_1=(0,0),&p_2=(1,0),&p_3=(0,1),&p_4=(1,1),\\
p_5=(2,15),&p_6=(15,4),&p_7=(11,15),&p_8=(12,16).
\end{array}
\]
これらが定義3.4の一般の位置にあることを確認し，blow up $Y_p$ を考える．例外曲線を $e_i$ とし，line の類を
\[
[a;\mu_1,\ldots,\mu_8]:=ah-\sum_i\mu_i e_i
\]
と表す．$a>0$ なら $\mu_i$ は平面上の像の $p_i$ における重複度であり，例外曲線そのものには負の係数が現れる．必要な9本は以下の通りである．
\[
\begin{array}{c|r|rrrrrrrr}
 &a&\mu_1&\mu_2&\mu_3&\mu_4&\mu_5&\mu_6&\mu_7&\mu_8\\\hline
\ell_1&0&-1&0&0&0&0&0&0&0\\
\ell_2&3&2&1&1&1&1&1&1&0\\
\ell_3&6&3&2&2&2&2&2&2&2\\
\ell_4&1&1&1&0&0&0&0&0&0\\
\ell_5&2&1&0&1&1&1&1&0&0\\
\ell_6&3&2&0&1&1&1&1&1&1\\
\ell_7&2&1&1&1&1&1&0&0&0\\
\ell_8&6&2&3&2&2&2&2&2&2\\
\ell_9&6&2&2&3&2&2&2&2&2
\end{array}
\]
平面曲線の方程式は，次数と指定した重複度を満たす係数の線形方程式を解いて求める．交点数は $aa'-\sum_i\mu_i\mu_i'$ で計算できるが，横断性と共通点の有無には実際の方程式の確認が必要である．論文では次の組を使う．
\[
\begin{array}{c|c@{\qquad}c|c}
\text{二つ組}&m&\text{三つ組}&t\\\hline
\ell_1,\ell_2&2&\ell_1,\ell_4,\ell_5&[1,1,1]\\
\ell_1,\ell_3&3&\ell_1,\ell_4,\ell_6&[1,1,2]\\
&&\ell_1,\ell_3,\ell_7&[1,1,3]\\
&&\ell_1,\ell_6,\ell_9&[1,2,2]\\
&&\ell_1,\ell_6,\ell_8&[2,2,2]
\end{array}
\]
すべての組に $\ell_1=e_1$ が入る点が便利である．$e_1\cong\PP^1$ と狭義変換の交点は，平面曲線の $p_1$ での接方向である．二つ組では，$\ell_2$ の二次接錐，$\ell_3$ の三次接錐がそれぞれ相異なる根を持つことを確認する．三つ組では，他の2曲線の接錐が共通の根を持たないことを確認する．根は $\overline{\mathbb F}_{19}$ 上で考える．これにより所要の横断性と三重点の不存在が得られる．

有限体上の例で十分なのは，これらの条件が代数的な開条件だからである．8点の整数持ち上げを取り，重複度を指定する線形系の階数条件と，接錐の判別式・終結式の非零性を保つ．19を法として非零の量は標数0でも非零なので，複素数上にも同じ性質の例が得られる．論文の有限体計算をこのように標数0へ移し，前頁の既約性と開性を使えば命題4.4が従う．
\par\smallskip
\noindent{\footnotesize 参照：I. Shimada, \emph{Del Pezzo surfaces of degree one and examples of Zariski multiples}，\S3，\S4.2--4.4．群構造 $2.G.2$ は Proposition 3.8 の証明，有限体上の例は Proposition 4.4 の証明（指定PDF pp.\,13--14）による．方程式の具体的な係数は論文の計算データに譲る．}
\end{document}
